\Needspace{12\baselineskip}
\subsubsection{3.2.12 UC12 登记装备故障并隔离库存}\label{uc-v3-12}
\begin{BSRUsecase}{UC12 登记装备故障并隔离库存}
主执行者 & 店员、门店经理 \\
辅执行者 & 无独立外部执行者；数据库、SSE 与计时模块均为系统内部机制。 \\
涉众利益 & 顾客需要安全装备；店员需要立即停止故障装备出租；技师需要明确的问题、负责人和截止时间；经理需要控制停机影响。 \\
目标与触发 & 发现可用装备故障，或归还验收发现损坏，创建可追溯维修任务。 \\
前置条件 & 操作者具有相应权限；手工登记单件为 AVAILABLE；不存在同一单件未关闭工单。归还损坏由归还流程在同一事务中自动生成。 \\
成功后置条件 & 生成 ASSIGNED 工单及首条时间线；单件为 MAINTENANCE；库存流水、审计和站内通知可查询。 \\
最小保证 & 失败不提交部分业务变更；权限或版本冲突返回明确原因。业务事实以服务器记录为准。 \\
基本流程 & \begin{minipage}[t]{\linewidth}\raggedright 1 执行者提交库存编号、故障说明、技师、优先级与预计完成日期。\\ 2 系统校验单件状态、开放工单、负责人和输入范围。\\ 3 系统在事务中创建工单，将单件隔离并写入库存流水。\\ 4 系统记录创建者、时间、初始状态、版本 1 和故障原文。\\ 5 系统保存站内通知；成功提交后发布变更信号。\\ 6 系统返回工单编号和隔离结果；获分配技师可在看板查询。\strut\end{minipage} \\
扩展流程 & \begin{minipage}[t]{\linewidth}\raggedright 2a 单件已借出、在途、维修或报废：手工登记被拒绝，原库存保持不变。\\ 2b 已有开放工单：提示避免重复登记，不另建第二张开放工单。\\ 2c 描述少于 3 字或多于 1000 字、负责人非技师、优先级或日期无效：拒绝。\\ 3a 事务失败：工单与库存更新全部回滚。\strut\end{minipage} \\
字段与状态事件 & unitId、issue、contractorId、priority、dueAt；服务器生成 id、storeId、createdAt、updatedAt、status、version；事件为故障登记、库存隔离和工单创建。 \\
\rowcolor{GoldPale}业务规则 & BR-035、BR-039、BR-044；一工单对应一物理单件和当前一负责人；优先级 LOW/NORMAL/HIGH/URGENT；未完成工单不能重复登记。 \\
质量与验收 & 故障创建成功后再次查库存为维修隔离；技师不能看其他人任务；计时目标为联网空闲页面 5 秒内刷新，须专项测试。 \\
设计约束 & 手工登记使用 POST /api/maintenance；事务内写数据库；站内通知已实现，外部短信邮件和可靠重试队列仍为建设目标。 \\
\end{BSRUsecase}
